\section{The integral counts and the relative deficit}
\label{sec:counts}

Let $X_m$ be the complex nine-dimensional manifold constructed in
Proposition~\ref{geo:family}, and let $Y$ be the surface of
Proposition~\ref{dyn:surface}. Set
\[
 C_m=X_m\times\CP^1,\qquad M_m=X_m\times Y.
\]
Both products are simply connected and Kähler, of real dimensions $20$
and $22$. Our task is to compare the ordinary Morse count on $C_m$ with
the stable Morse count on $M_m$. By Theorem~\ref{morse:equality}, this
amounts to computing their quantities $\lambda$.

\subsection{The two torsion profiles}
The ranks in Lemma~\ref{geo:surfaces} give $b(S_2)=b(S_3)=8$.
Since $\CP^1\times\CP^1$ has free homology of ranks $(1,2,1)$ in degrees
$(0,2,4)$, the center $Z_3=S_3\times\CP^1\times\CP^1$ has
$b(Z_3)=32$. The integral blow-up decomposition in
\eqref{geo:family-homology} therefore yields
\begin{equation}\label{counts:rank-X}
 b(X_m)=10+m(6\cdot8+4\cdot32)=10+176m.
\end{equation}

All the homology shifts in that decomposition are even. Torsion in the
surface degrees $1$ and $2$ consequently gives two separate lists, in
degrees $1+2j$ and $2+2j$. For either list we have
\begin{equation}\label{counts:base-profiles}
 \Tor H_{1+2j}(X_m;\Z)\cong\Tor H_{2+2j}(X_m;\Z)
 \cong(\Z/2)^{ma_j}\oplus(\Z/3)^{mq_j},
\end{equation}
where the finite sequences, indexed starting at $j=1$, are
\[
 a=(2,2,2,2,2,2),\qquad q=(1,3,4,4,3,1),
\]
and all other entries are zero. To see the second sequence directly, let
$u$ record a shift by two homological degrees. The product with
$\CP^1\times\CP^1$ and the four blow-up shifts give
\[
 (u+u^2+u^3+u^4)(1+2u+u^2)
 =u+3u^2+4u^3+4u^4+3u^5+u^6.
\]
The first sequence records six blow-up shifts of two copies of $\Z/2$.
Equation~\eqref{counts:base-profiles} accounts for all torsion.

For nonnegative integers $a,q$, the primary-generator identity
\eqref{morse:primary-maximum} gives
\begin{equation}\label{counts:generator-maximum}
 d\bigl((\Z/2)^a\oplus(\Z/3)^q\bigr)=\max(a,q).
\end{equation}
In particular, generator counts must be computed in each degree
before summing over degrees.

If a factor $B$ has torsion-free homology concentrated in even degrees, put
$e_k=\rank H_{2k}(B;\Z)$ for $k\geq0$. The integral Künneth theorem replaces $a$ and $q$
by their convolutions with $e$:
\[
 (e*a)_j=\sum_{k\geq0}e_k a_{j-k},\qquad
 (e*q)_j=\sum_{k\geq0}e_k q_{j-k}.
\]
There are no additional Künneth torsion terms. For $C_m$ use $e=(1,1)$;
for $M_m$, \eqref{dyn:Y-homology} gives $e=(1,6,1)$. The resulting
profiles in either degree list, divided by $m$, are shown in
Table~\ref{tab:profiles}.

\begin{table}[tb]
\centering
\small
\begin{tabular}{@{}llr@{}}
\toprule
Product and primary part & Multiplicities in one degree list & Sum\\
\midrule
$C_m$, $2$-primary & $(2,4,4,4,4,4,2)$ & $24$\\
$C_m$, $3$-primary & $(1,4,7,8,7,4,1)$ & $32$\\
$C_m$, generators & $(2,4,7,8,7,4,2)$ & $34$\\
\midrule
$M_m$, $2$-primary & $(2,14,16,16,16,16,14,2)$ & $96$\\
$M_m$, $3$-primary & $(1,9,23,31,31,23,9,1)$ & $128$\\
$M_m$, generators & $(2,14,23,31,31,23,14,2)$ & $140$\\
\bottomrule
\end{tabular}
\caption{Primary torsion multiplicities and their entrywise maxima.
Every displayed entry is multiplied by $m$, and the same list occurs
once in odd degrees and once in even degrees. The generator row takes
the maximum within each column, as in
\eqref{counts:generator-maximum}.}
\label{tab:profiles}
\end{table}

It follows that
\begin{align}
 \sum_i t_i(C_m)&=2\cdot34m=68m,&
 \sum_i t_i(M_m)&=2\cdot140m=280m.\label{counts:torsion-totals}
\end{align}
The free ranks are $b(C_m)=2b(X_m)$ and $b(M_m)=8b(X_m)$, so
\begin{equation}\label{counts:lambda}
 \lambda(C_m)=20+488m,\qquad
 \lambda(M_m)=80+1968m.
\end{equation}
The table also explains the sign of the gap. The total multiplicity at
each individual prime is four times as large for $M_m$ as for $C_m$.
However, relative to the $3$-primary row, taking the degreewise maximum
adds $2m$ generators per parity for $C_m$ and $12m$ for $M_m$. Thus the
generator count for $M_m$ exceeds four times that for $C_m$ by
$(12-4\cdot2)m=4m$ per parity. The free-rank contributions cancel,
and the two parities and the factor $2$ in $\lambda$ give
\begin{equation}\label{counts:gap}
 \lambda(M_m)-4\lambda(C_m)
 =2(280-4\cdot68)m=16m.
\end{equation}

\subsection{Completion of the construction}
\begin{proof}[Proof of Theorem~\ref{thm:main}]
Fix $m\geq1$. Choose any Kähler form on $X_m$ and equip
$M_m=X_m\times Y$ with its product with the invariant form on $Y$ from
Proposition~\ref{dyn:surface}. Theorem~\ref{morse:equality} applies to
$C_m$ and $M_m$, since both are closed and simply connected and their
real dimensions are at least $9$. It supplies an ordinary Morse function
$h_m$ on $C_m$ with $20+488m$ critical points, and it gives
$\SM(M_m)=80+1968m$.

Corollary~\ref{dyn:product}, applied to $X_m$ and $h_m$, produces a smooth
one-periodic Hamiltonian $H_m$ on $M_m$ with precisely
\[
 4\#\Crit(h_m)=80+1952m
\]
fixed points. They are all nondegenerate. Simple connectivity of $M_m$
makes every corresponding loop along $\phi_{H_m}^t$ contractible, so the
two fixed-point counts agree.

Finally,
\[
 \frac{\SM(M_m)-\#\Fix(\phi_{H_m}^1)}{\SM(M_m)}
 =\frac{16m}{80+1968m}
 =\frac{m}{5+123m}\geq\frac1{128},
\]
where the inequality is equivalent to $m\geq1$.
For a given $R>0$, take
$m=\max\{1,\lceil(R-80)/1968\rceil\}$. Then
$\SM(M_m)\geq R$, while the dimension and $\delta=1/128$ remain fixed.
\end{proof}

\begin{remark}\label{counts:field-bounds}
The total Betti numbers of $M_m$ over fields of characteristic $2$ and $3$
are respectively $80+1792m$ and $80+1920m$. Indeed, the universal coefficient
theorem adds twice the total primary multiplicity to the integral free
rank, and Table~\ref{tab:profiles} gives those multiplicities. Both
numbers are smaller than the fixed-point count $80+1952m$. Over any
field of characteristic different from $2$ and $3$, the total Betti
number is $80+1408m$. Thus taking a single
coefficient field does not recover the integral generator count used here:
the prime contributing the larger number of generators changes with the
degree.
\end{remark}

\begin{remark}\label{counts:integral-Floer-bound}
The integral Floer bound of Bai--Xu uses homology graded modulo twice the
minimal Chern number~\cite[Theorem~A]{BaiXu2022}; see also
\cite[Corollary~1.2]{BaiXu2026}. The minimal Chern number $N$ is the positive generator of
the nonzero subgroup $c_1(TM_m)(\pi_2(M_m))\subset\Z$.
For $M_m$ it is $1$: if $E\subset Y$ is an exceptional line, then its
tangent bundle has degree $2$ and its normal bundle has degree $-1$.
Thus $c_1(TM_m)$ evaluates to $1$ on the sphere $\{x\}\times E\subset M_m$.
In the grading modulo $2$, all odd-degree groups are combined and all
even-degree groups are combined. In either parity, Table~\ref{tab:profiles}
then gives primary torsion multiplicities $96m$ and $128m$. The minimum
number of torsion generators is therefore $128m$ in each parity. The
cyclically graded integral bound is
\[
 b(M_m)+2(128m+128m)=80+1920m,
\]
which is smaller than the fixed-point count by $32m$. The strict gap in
Theorem~\ref{thm:main} relies on keeping the integer homological degrees
separate when computing the stable Morse number.
\end{remark}
